\section{摘要}

真实量子芯片通常只支持固定连接上的双比特门，而理想量子电路往往默认任意
逻辑量子比特都可以直接交互。量子路由编译的任务是在保持电路语义等价的前提
下，把逻辑量子比特映射到物理拓扑上，并插入必要的 SWAP 门，使电路能够在目标
硬件上合法执行。该问题可以表述为图约束下的映射与搜索问题，具有明显的组合
优化特征。

本项目实现了一条神经辅助量子路由流程 NPQR。该流程以神经网络动作为核心评分
信号，同时结合初始映射选择、有界束搜索、前沿触发、动作剪枝、局部后缀修复和
轨迹复放验证。SABRE 作为稳定的传统启发式基线，仅用于结果对比，不参与 NPQR
的路线生成。系统提供命令行、REST API、MCP 服务和网页演示入口，使算法能够被
直接调用和展示。

基于现有公开证据，项目已经完成可运行、可验证、可解释的组合式量子路由流程。
快速基线矩阵显示 SABRE 在不同启发式配置下具有稳定表现；项目证据同时表明
NPQR 能在固定困难样例上完成路由、返回可复放轨迹，并保持 REST API 与 MCP 调用
结果一致。本文报告算法建模、流程设计、系统实现、复杂度分析和结果边界。报告
不声称 NPQR 在所有电路上优于 SABRE，也不声称默认模型达到最先进量子路由水平。
